\subsection{Mini-Batch Attenuation of Directional Alignment}
\label{sec:sgd-noise-experiments}

To test the effect of sampling, we vary batch size at initialization
and after 2,000 GD updates in the same three model--dataset cases.
Learning rates remain fixed at $0.02$, $0.1$, and $0.2$ for the MLP,
ResNet, and Transformer, respectively. Each starting state receives
200-update continuations, with three initializations, eight sampling
repetitions per stochastic condition, and one full-GD control per
initialization. Figure~\ref{fig:sgd-moving-geometry} summarizes mean
consecutive-gradient cosine over these continuations. Sampling units
and objective coverage are specified in
Appendix~\ref{appendix:gd-raw-alignment}; equal update budgets imply
different amounts of data exposure.

From the trained states, the smallest batches move the strongly
negative full-GD alignment toward zero in all three cases. This
attenuation is consistent with sampling weakening the coherent
back-and-forth motion of the destructive regime. The response need
not be monotone across intermediate batch sizes: sampling changes
both the measured gradient and the trajectory itself.

The initial-state comparison makes this dependence on the trajectory
visible. Full GD has positive mean cosine in the MLP and ResNet,
but negative mean cosine in the Transformer, where reversal develops
within the 200-update window. Smaller MLP batches can also induce
negative mean alignment despite a positive full-GD mean. Thus the
batch size influences which regime training reaches as well as the
coherence within that regime.

Directional attenuation alone does not determine objective progress.
For the Transformer at $b=8$, mean cosine is near $-0.26$ from both
starting states, yet mean full-data loss falls
from initialization and rises from the trained state. Together with
the full-batch results, this motivates reading alignment alongside
gradient magnitude and loss: the same directional regime can support
descent or an upward excursion depending on the gradient's growth
and the state being traversed.
